\chapter{150 MHz에 연산을 담기까지}

\section{성공한 회로만 보면 설계 이유가 사라진다}

최종 RTL만 읽으면 register가 왜 그 위치에 있는지 알기 어렵다. 이 설계에서는
clock constraint를 바꾸고 top timing path를 다시 읽을 때마다 병목의 종류가
달라졌다. 남아 있는 report와 RTL milestone을 시간 순서로 놓으면 다음과 같다.

\begin{longtable}{L{22mm}L{50mm}L{68mm}}
\toprule
시점 & 관찰한 경로와 결과 & 다음 판단 \\
\midrule
최초 100~MHz &
DUMP address $\rightarrow$ BRAM address,
9.077 ns &
주소, bank, phase 선택을 BRAM pin 앞에서 끝내는 구조는 더 높은 주파수에
맞지 않았다. \\
초기 150~MHz &
mode select $\rightarrow$ radix-3 subtraction register,
$\wns=-0.338$ ns &
한 cycle에 연결된 23-bit carry 연산과 뒤쪽 선택을 분리해야 했다. \\
200~MHz 산술 실험 &
DSP product $\rightarrow$ ML-DSA intermediate register,
8.168 ns &
product 결합과 reduction digit 계산을 한 register 경계 안에 두지 않기로 했다. \\
최종 150~MHz &
standalone, AXI4-Stream, BRAM complete system 모두 timing 충족 &
목표 주파수를 검증된 production point로 고정하고 세 범위를 따로 보존했다. \\
\bottomrule
\end{longtable}

이 표의 각 행은 동일한 단일 변수 실험이 아니다. 개발 중 서로 다른 snapshot을
나열한 설계 일지다. 따라서 수치 차이 전체를 특정 RTL 변경의 효과로 해석하지
않는다. 여기서 사용할 수 있는 증거는 “그 시점의 top path가 무엇이었는가”와
“그 path의 조합 구조를 다음 RTL에서 어떻게 끊었는가”다.

\section{몇 ns가 걸리는지에 대한 직관}

RTL 연산 이름만으로 지연을 고정값처럼 외우면 안 된다. 같은 23-bit addition도
앞의 operand MUX, 뒤의 result MUX, fanout과 배치 위치에 따라 route delay가
달라진다. 이 프로젝트에서 실제 관찰한 path를 구성별로 나누면 다음과 같다.

\begin{longtable}{L{45mm}r r r L{45mm}}
\toprule
실제 경로 & Logic & Route & Total & 관찰한 구조 \\
\midrule
최초 DUMP address $\rightarrow$ BRAM &
1.681 ns & 7.396 ns & 9.077 ns &
8 LUT, route 81.5\% \\
초기 radix-3 subtraction result &
2.896 ns & 3.781 ns & 6.677 ns &
11 CARRY4를 포함한 13 levels \\
200~MHz ML-DSA intermediate &
4.600 ns & 3.568 ns & 8.168 ns &
12 CARRY4를 포함한 17 levels \\
최종 standalone modular-add register &
2.494 ns & 4.028 ns & 6.522 ns &
4 CARRY4를 포함한 7 levels \\
최종 BRAM system $1006x$ table-sum &
2.920 ns & 3.630 ns & 6.550 ns &
6 CARRY4를 포함한 10 levels \\
\bottomrule
\end{longtable}

여기서 얻는 실무적 직관은 세 가지다.

\begin{enumerate}
  \item 23-bit carry 연산 두 개와 mode selection이 이어진 path는 이
        implementation에서 6.677~ns였고 150~MHz에 들어가지 않았다.
  \item 더 복잡한 product 결합과 field 계산은 8.168~ns였다. 5~ns에
        맞추려면 식 정리만으로 부족하고 register 또는 DSP 내부 stage가 필요했다.
  \item logic이 짧아도 BRAM 주변처럼 route가 80\%를 넘으면 9~ns가 될 수 있다.
        반대로 “adder는 몇 ns”라는 표만으로 system timing을 예측할 수 없다.
\end{enumerate}

이 수치는 xc7z020 계열, Vivado 2020.2, 해당 placement와 constraint에서 얻은
경로 측정값이다. 다른 device, speed grade, floorplan과 fanout에는 그대로
대입하지 않는다. 설계 초기에 필요한 것은 절대값 암기가 아니라
“carry가 몇 번 직렬인가”, “MUX가 carry 앞뒤 어디에 있는가”,
“logic보다 route가 큰가”를 구분하는 감각이다.

\section{실패 1: subtract와 correction을 한 cycle에 연결했다}

modular subtraction은 식으로 쓰면 짧다.
\[
d=a-b,\qquad r=
\begin{cases}
d+q,&d<0,\\
d,&d\ge0.
\end{cases}
\]
그러나 23-bit RTL에서 raw subtraction, borrow 선택, conditional addition과
mode-dependent result 선택이 같은 register 앞에 놓이면 carry chain 두 개와
MUX가 직렬로 합성될 수 있다.

초기 150~MHz post-route report의 top path는 mode register에서
radix-3 subtraction result register까지였다. 6.667~ns requirement에서 data
path가 6.677~ns였고, logic 2.896~ns와 route 3.781~ns를 사용했다. 13 logic
levels 가운데 CARRY4가 11개였다. setup slack은 $-0.338$~ns였다. 이 결과는
“subtractor 하나가 느리다”보다 다음 사실을 보여 준다.

\begin{itemize}
  \item carry 결과를 고르는 mode MUX도 같은 path의 일부다.
  \item 산술 depth와 route가 모두 커서 placement directive만으로 해결하기 어렵다.
  \item source가 mode register라는 사실은 연산뿐 아니라 control fanout도
        함께 줄여야 함을 뜻한다.
\end{itemize}

최종 RTL은 모든 modular subtraction에 새 stage를 추가하지 않았다. multiplier
앞에서 일찍 필요한 두 general difference만 C1과 C2로 나눴다.

\begin{lstlisting}
// C1: raw subtraction and borrow
raw0_c1    <= {1'b0, b} - {1'b0, a};
borrow0_c1 <= raw0[D];

// C2: conditional correction
corrected = {1'b0, raw0_c1[D-1:0]} + {1'b0, q};
d0_c2     <= borrow0_c1 ? corrected[D-1:0]
                        : raw0_c1[D-1:0];
\end{lstlisting}

C2는 이 rail에서 원래 alignment에 쓰이던 자리였다. 그러므로 새 latency를
더하지 않고 기존 register의 일을 바꾸어 두 carry chain을 서로 다른 cycle에
배치했다. 이 사례에서 핵심 판단은 “pipeline을 한 단계 늘린다”가 아니라
\textbf{이미 존재하는 alignment stage에 실제 연산을 배치한다}는 것이다.

또한 arithmetic result인 radix-3의 $E$와 $S$를 첫 register에서 공통 A/C rail에
바로 합치지 않았다. carry chain 직후에 mode MUX가 붙기 때문이다. 첫 register는
분리하고, 다음 경계에서 이미 등록된 값끼리 merge했다.

\section{실패 2: ML-DSA product 결합과 digit 계산을 함께 수행했다}

200~MHz 실험의 한 post-route top path는 DSP48E1의 product output에서
\rtlfile{a\_tmp\_buf} register까지였다. 5.000~ns requirement에 data path는
8.168~ns, logic levels는 17개였고 그중 CARRY4가 12개였다. slack은
$-3.171$~ns였다.

당시 조합 cone은 low product의 정렬된 high part와 다른 partial product를
더한 뒤, 그 결과에서 여러 field sum을 연속으로 만들었다. RTL의 수식은 맞고
simulation도 통과했지만, Vivado가 본 물리 경로에는 긴 carry chain이 남았다.
“중간 wire 이름을 여러 개 둔다”는 것은 pipeline이 아니다. register edge가
없으면 모든 식은 여전히 한 cycle의 경로다.

최종 six-stage \rtlfile{unifiedMUL}은 일을 다음처럼 나눈다.

\begin{enumerate}
  \item DSP0가 16-bit low product를 계산한다.
  \item DSP1의 MREG/PREG가 ML-DSA high partial product와 aligned cross-sum을
        받아 외부 fabric의 긴 product-adder 경계를 줄인다.
  \item C3에서 $H_0+H_1$을 한 번 계산하고 다음 high-part 식에서도 재사용한다.
        별도 24-bit addition 대신 두 번째 carry chain을 14 bit로 제한한다.
  \item C4에서 segmented residual, C5에서 base/correction addition,
        C6에서 constant-selected canonical correction을 수행한다.
\end{enumerate}

마지막 correction도 high digit과 low digit을 각각 고치는 branch로 표현했을 때
합성기가 branch MUX와 carry를 깊게 합쳤다. 최종판은 범위 판정으로
$0,+q,-q,-2q$ 중 correction constant 하나를 고르고 signed addition 한 번만
수행한다. 수학적으로 같은 식 가운데 FPGA에서 carry와 MUX가 직렬로 놓이지
않는 표현을 선택한 것이다.

\section{실패 3: Montgomery의 다섯 operand를 한 번에 더했다}

초기 R16 Montgomery path는 선택된 네 shifted operand를 더해
$c_qT\bmod 2^9$를 만들고, 같은 조합 경로에서 $T[15:7]$까지 더해
$u[15:7]$을 완성했다.

\begin{lstlisting}
shift_sum = op0 + op1 + op2 + op3;
u_hi      = shift_sum + T[15:7];
\end{lstlisting}

operand 폭은 9 bit로 작지만 다섯 항의 mode selection과 addition이 한 cycle에
모였다. 최종판은 C2에 $c_qT\bmod512$만 등록하고, 등록된 $T[15:7]$은 다음
DSP-input cycle에서 더한다.

\begin{lstlisting}
// C2 register
corr9_c2 <= op0 + op1 + op2 + op3;
t_c2     <= T;

// following DSP-input cycle
u_hi = t_c2[15:7] + corr9_c2;
u    = {u_hi, t_c2[6:0]};
\end{lstlisting}

DSP1이 ML-DSA와 Montgomery mode에서 서로 다른 cycle 역할을 갖는다는 schedule을
이용했으므로 six-cycle latency와 \ii=1을 유지했다. pipeline 최적화는 연산을
기계적으로 반으로 자르는 일이 아니라, 다음 hard block이 operand를 받는
시점까지 포함해 cut 위치를 정하는 일이다.

\section{실패 4: 상수곱에 stage 하나를 더 넣자 기능이 깨졌다}

NTRU+의 고정 상수곱은 처음에 shifted partial product 일곱 개와
$q=4096-639$를 이용한 여러 fold를 두 cycle에 배치했다. adder tree를
balanced form으로 다시 써도 product와 modular fold의 carry depth 자체는
남았다. 이후 작은 input field를 table로 바꾸어 carry chain을 table output의
합과 한 번의 correction으로 줄였다. 최종 $1006x\bmod3457$ 회로는
6-bit field 두 개의 table과 12-bit modular addition을 사용한다.

더 높은 주파수 실험에서는 이 상수곱을 세 stage로 늘려 보았다. 개별 block의
조합 depth는 줄지만 radix-3/삼항식 branch는 상수곱 결과와 별도의 difference
rail이 같은 cycle에 도착한다고 가정했다. 상수곱만 한 cycle 늦어지자 NTRU+
INTT 결과가 어긋났다.

선택지는 두 가지였다.

\begin{enumerate}
  \item 관련된 모든 branch와 metadata를 한 cycle 더 지연한다.
  \item 두-cycle 계약을 유지하고 연산 표현을 바꾼다.
\end{enumerate}

첫 선택은 공통 alignment FF와 전체 latency를 늘린다. 최종판은 두 번째를
택했다. raw subtraction의 underflow correction을 table 앞 12-bit adder로
수행하지 않고, underflow일 때 더해야 하는 상수 효과
$-1006\times639\equiv168\pmod{3457}$를 table 선택에 미리 반영했다.
그 결과 correction carry chain 하나를 table data selection으로 바꾸면서
두-cycle arrival contract를 보존했다.

이 실패는 “register를 넣으면 timing은 좋아진다”는 문장이 불완전함을
보여 준다. multi-branch pipeline에서는 값, valid, mode, destination tag와
parallel operand가 같은 transaction으로 다시 만나는 cycle까지 함께 바뀐다.

\section{실패 5: 연산기를 최대한 공유했더니 더 커졌다}

한 실험판은 네 개의 runtime-generic folded ALU에 23-bit와 12-bit 연산을 모두
모으고, raw/correction two-pass register와 공통 delay rail을 두었다. 연산기
개수만 세면 매력적이지만 실제 FPGA 비용은 다음 위치로 이동했다.

\begin{itemize}
  \item generic ALU 앞의 wide operand crossbar
  \item carry chain 뒤의 mode-dependent result MUX
  \item 여러 stage와 mode로 퍼지는 select fanout
  \item latency가 다른 branch를 끝까지 끌고 가는 alignment FF
\end{itemize}

NTRU+ lazy reduction도 다섯 internal correction을 제거하는 대신 intermediate가
13--14 bit로 넓어졌다. multiplier와 constant-multiplier 입력이 커지고 output
canonicalization이 다시 필요했다. 최종 output에서 제거했던 adder/MUX가
돌아오므로 해당 branch의 순 area 이득이 사실상 사라졌다.

최종 구조는 큰 자원과 시점이 정확히 맞는 연산을 공유한다. 두 DSP, DSP
register, final correction addition은 공유한다. NTRU+의 cycle-local add/sub는
12 bit로 남긴다. 이 경험에서 나온 기준은 다음과 같다.

\begin{tradeoffbox}[공유 여부를 결정하는 기준]
연산기 개수만 비교하지 않는다. 공유 전후의 operand MUX 폭, result MUX 위치,
mode fanout, alignment register와 placement 거리까지 합쳐서 비교한다.
\end{tradeoffbox}

\section{implementation directive만으로 150 MHz가 되지 않았다}

중간 RTL을 Zynq-7020에 standalone OOC로 구현했을 때 default flow의
150~MHz WNS는 $-0.374$~ns였다. synthesis retiming과 timing-oriented
place/route directive를 적용하면 $-0.152$~ns까지 회복했지만 여전히
timing은 닫히지 않았다. 같은 flow에서 140~MHz는 $+0.026$~ns였다.

이 수치는 directive가 쓸모없다는 뜻이 아니다. RTL의 register, MUX, fanout
구조를 고친 뒤 implementation flow가 나머지 placement/routing margin을
회복한다. 최종 150~MHz 결과는 두 작업을 함께 적용한 결과다.

300~MHz 실험에서는 3.333~ns constraint에서 synthesis WNS $-7.170$~ns와
RAMB36E1 pulse-width slack $-0.551$~ns가 관찰되었다. 이 implementation
point는 route를 진행하지 않았다. pulse-width violation은 조합 path에 register를
더 넣는 것으로 해결되지 않기 때문이다. 목표 주파수를 올리는 실험도
\textbf{RTL로 고칠 path}, \textbf{clock/device primitive 제한},
\textbf{기능 정렬 비용}을 구분해 중단 기준을 세워야 한다.

\begin{checklistbox}
\begin{itemize}
  \item 실패한 run도 requirement, WNS, source, destination, logic levels,
        logic/route 비율을 남긴다.
  \item 한 cycle 식에서 carry chain이 몇 번 이어지는지 회로로 다시 그린다.
  \item 빈 alignment stage를 먼저 찾고 새 stage 추가는 그다음에 검토한다.
  \item 한 branch의 latency를 바꾸면 parallel data와 metadata의 합류 cycle을
        testbench에서 확인한다.
  \item 공유 실험은 ALU 수뿐 아니라 MUX, fanout, FF와 route까지 비교한다.
  \item directive로 얻은 margin과 RTL 구조 변경의 효과를 같은 것으로 쓰지 않는다.
\end{itemize}
\end{checklistbox}
